\section{Method}
\label{sec:method}

\begin{figure*}[t]
  \centering
  \makebox[\linewidth][c]{% Overview figure (figure*, drawn at \textwidth = 17.46 cm). The Overleaf figure (fig/overview_v1.tex), refined:
% numbered stages with a one-line description each, single-turn inference only (no KV cache), less empty space,
% larger frames and text. Three stages left to right:
% (1) retrieval: video + subtitles -> retriever (with query q) -> relevance r_i over items in time order;
% (2) Token Scheduler (the contribution, the only accented stage): objective, value curves v_frame / v_text,
%     water-filling -> token plan b_i;
% (3) inference: frozen pruner keeps each frame's top-ranked tokens, the frozen VLM answers.
% Bar heights and the plan are illustrative, not measured. Frames: Video-MME video bepwr1-CNRU (illustration only).
% Tableau 10 (Stone, 2016): blue 4E79A7 and orange F28E2B on neutral greys; dark text on light fills only
\definecolor{ink}{HTML}{333333}%
\definecolor{arrowgray}{HTML}{6F6F6F}%
\definecolor{bandA}{HTML}{F4F5F7}%
\definecolor{bandB}{HTML}{EAF1F8}%
\definecolor{edgeB}{HTML}{4E79A7}%
\definecolor{bandC}{HTML}{F4F5F7}%
\definecolor{stepA}{HTML}{333333}%
\definecolor{stepB}{HTML}{4E79A7}%
\definecolor{stepC}{HTML}{333333}%
\definecolor{modelA}{HTML}{DCE6F1}%
\definecolor{modelC}{HTML}{DCE6F1}%
\definecolor{modelD}{HTML}{4E79A7}%
\definecolor{modelT}{HTML}{333333}%
\definecolor{flake}{HTML}{4E79A7}%
\definecolor{chipF}{HTML}{FFFFFF}%
\definecolor{chipD}{HTML}{A0A0A0}%
\definecolor{frm}{HTML}{4E79A7}%
\definecolor{txt}{HTML}{F28E2B}%
\definecolor{accA}{HTML}{4E79A7}%
% every color role: ink, arrowgray, bandA/B/C, edgeB, stepA/B/C, modelA/C/D/T, flake, chipF/D, frm, txt
\newcommand{\snowflake}[2]{%
  \begin{scope}[shift={#1}, line width=0.55pt, draw=flake, line cap=round]
    \foreach \a in {0,60,...,300}{\draw (0,0) -- (\a:#2);
      \draw (\a:0.55*#2) -- ++(\a+45:0.32*#2); \draw (\a:0.55*#2) -- ++(\a-45:0.32*#2);}
  \end{scope}}
% \items{width}{height scale}{v1/t1,...,v20/t20}: 20 items on one time axis; t = 1 marks a transcript event (light bar)
\newcommand{\items}[3]{%
  \foreach \v/\t [count=\i] in {#3} {
    \ifnum\t=1 \fill[txt] ({(\i-0.5)*#1/20-0.3*#1/20},0) rectangle ++({0.6*#1/20},{\v*#2});
    \else      \fill[frm]    ({(\i-0.5)*#1/20-0.3*#1/20},0) rectangle ++({0.6*#1/20},{\v*#2}); \fi}}
% \prunedframe{center}{image}{kept tokens out of 32}: a frame (1.0 cm) as an 8x4 token grid; dropped tokens greyed out.
% The keep order is one fixed ranking (centre first), so a frame at b tokens holds its top b tokens.
\newcommand{\prunedframe}[3]{%
  \node[inner sep=0pt] at (#1) {\includegraphics[width=1.0cm]{fig/frames/#2}};
  \begin{scope}[shift={($(#1)+(-0.5,-0.28125)$)}]
    \foreach \c [count=\k from 0] in {10,11,18,19,9,12,17,20,2,3,26,27,1,4,25,28,8,13,16,21,0,5,24,29,6,14,22,30,7,15,23,31} {
      \ifnum\k<#3\relax\else
        \pgfmathtruncatemacro{\cx}{mod(\c,8)}\pgfmathtruncatemacro{\cy}{int(\c/8)}
        \fill[white, opacity=0.85] (\cx*0.125,\cy*0.140625) rectangle ++(0.125,0.140625);
      \fi}
    \draw[ink!30, line width=0.15pt, xstep=0.125, ystep=0.140625] (0,0) grid (1.0,0.5625);
    \draw[ink!45, line width=0.3pt] (0,0) rectangle (1.0,0.5625);
  \end{scope}}
\newcommand{\txtchip}[1]{% transcript event: a short text chip
  \fill[txt!35, rounded corners=1pt] ($(#1)+(-0.42,-0.2)$) rectangle ++(0.84,0.4);
  \draw[txt!85!black, line width=0.5pt] ($(#1)+(-0.3,0.07)$) -- ++(0.6,0);
  \draw[txt!85!black, line width=0.5pt] ($(#1)+(-0.3,-0.07)$) -- ++(0.4,0);}
% \stage{x}{y}{number}{title}{one-line description}: numbered stage heading
\newcommand{\stage}[5]{%
  \node[circle, fill=#5, text=white, font=\sffamily\scriptsize\bfseries, inner sep=0pt, minimum size=0.4cm] at (#1+0.2,#2) {#3};
  \node[font=\sffamily\small\bfseries, text=#5, anchor=west, inner sep=0pt] at (#1+0.5,#2) {#4};}

\begin{tikzpicture}[font=\sffamily\footnotesize, text=ink, line width=0.6pt,
  card/.style={fill=white, draw=ink!18, line width=0.5pt, rounded corners=4pt, align=center, inner sep=3.5pt},
  arr/.style={-{Stealth[length=5pt, width=4pt]}, draw=arrowgray, line width=0.8pt},
  ax/.style={-{Stealth[length=3.5pt, width=3pt]}, draw=ink!55, line width=0.5pt},
  axl/.style={font=\sffamily\scriptsize, text=ink!70},
  desc/.style={font=\sffamily\scriptsize, text=ink!65, anchor=north west, inner sep=0pt, align=left},
  elbl/.style={font=\sffamily\scriptsize, text=ink!75, inner sep=1.5pt},
  model/.style={fill=modelC, draw=modelD, text=modelT, line width=0.5pt, rounded corners=3pt, align=center,
    minimum height=0.62cm, inner xsep=5pt, font=\sffamily\footnotesize},
  qa/.style={fill=chipF, draw=chipD, line width=0.5pt, rounded corners=3pt, align=left,
    inner sep=3pt, font=\sffamily\scriptsize}]

% ---------------- layout ----------------
% Top row: (1) Retrieval -> (2) Allocation; bottom row: (3) Inference, left to right. Stages connect box to box
% (r_i, b_i). Q sits inside Retrieval and A inside Inference; nothing is outside the stage boxes.
\fill[rounded corners=6pt, bandA] (0,2.0) rectangle (7.5,5.3);
\fill[rounded corners=6pt, bandB] (8.0,2.0) rectangle (17.4,5.3);
\draw[rounded corners=6pt, edgeB, line width=0.6pt] (8.0,2.0) rectangle (17.4,5.3);
\fill[rounded corners=6pt, bandC] (0,0) rectangle (17.4,1.6);
\stage{0.12}{5.02}{1}{Retrieval}{stepA}
\stage{8.12}{5.02}{2}{Allocation}{stepB}
\stage{0.12}{1.3}{3}{Inference}{stepC}
\draw[arr] (7.5,3.65) -- node[elbl, above] {$r_i$} (8.0,3.65);
\draw[arr] (12.7,2.0) -- node[elbl, right] {$b_i$} (12.7,1.6);

% ---------------- (1) retrieval ----------------
\fill[ink!88, rounded corners=1pt] (0.2,3.85) rectangle (7.3,4.75);
\foreach \x in {0.28,0.48,...,7.2} {\fill[white] (\x,4.65) rectangle ++(0.1,0.06); \fill[white] (\x,3.89) rectangle ++(0.1,0.06);}
\foreach \k in {1,...,5} {\node[inner sep=0pt] at ({0.91+(\k-1)*1.42},4.3) {\includegraphics[width=1.2cm]{fig/frames/frame\k.jpg}};}
\fill[txt!35, rounded corners=1pt] (0.2,3.57) rectangle (7.3,3.77);
\foreach \x/\l in {0.5/1.0, 1.9/0.8, 3.1/1.2, 4.7/0.9, 6.0/0.9} {\draw[txt!85!black, line width=0.5pt] (\x,3.67) -- ++(\l,0);}
\node[qa] (q) at (1.0,2.75) {\textbf{Q:} Color of\\the last car?};
\node[model, fill=modelA] (ret) at (2.95,2.75) {Retriever\hspace{0.12cm}};
\snowflake{($(ret.north east)+(-0.14,-0.14)$)}{0.075}
\draw[arr] (q.east) -- (ret.west);
\draw[arr] (2.95,3.55) -- (ret.north);
\begin{scope}[shift={(4.25,2.42)}]
  \draw[ax] (0,0) -- (3.05,0); \draw[ax] (0,0) -- (0,1.0);
  \node[axl, anchor=west] at (0.06,0.93) {$r_i$};
  \node[axl, anchor=north] at (1.5,-0.02) {frame \#};
  \items{2.95}{1.4}{.04/0,.06/0,.30/0,.52/0,.20/1,.40/0,.10/0,.05/0,.04/0,.36/1,
                   .06/0,.22/0,.30/0,.42/1,.34/0,.26/0,.08/0,.18/1,.05/0,.03/0}
\end{scope}
\draw[arr] (ret.east) -- (4.15,2.75);

% ---------------- (2) allocation ----------------
\node[font=\sffamily\scriptsize] at (12.7,4.5) {$\max\limits_{\mathbf b}\ \sum_i r_iv_i(b_i)-\eta\sum_{i\neq j}s_{ij}\,v_i(b_i)\,v_j(b_j)\qquad\text{s.t.}\ \ \sum_i b_i\le B$};
\begin{scope}[shift={(8.75,2.42)}]
  \draw[ax] (0,0) -- (3.2,0); \draw[ax] (0,0) -- (0,1.6);
  \node[axl, anchor=west] at (0.06,1.5) {$v(b)$};
  \node[axl, anchor=north] at (1.6,-0.02) {$b$};
  \draw[txt, line width=1.1pt, domain=0.02:8, samples=60] plot ({0.37*\x}, {1.45*((2^\x-1)^0.841)/(((2^\x-1)^0.841)+3.44)});
  \draw[frm, line width=1.1pt, domain=0.02:8, samples=60] plot ({0.37*\x}, {1.35*((2^\x-1)^0.841)/(((2^\x-1)^0.841)+7.73)});
  \node[font=\sffamily\scriptsize, text=txt!80!black, anchor=south west, inner sep=1pt] at (3.0,1.35) {$v_{\text{text}}$};
  \node[font=\sffamily\scriptsize, text=frm, anchor=north west, inner sep=1pt] at (3.0,1.2) {$v_{\text{frame}}$};
\end{scope}
\draw[arr] (12.75,3.2) -- (13.45,3.2);
\begin{scope}[shift={(13.85,2.42)}]
  \draw[ax] (0,0) -- (3.3,0); \draw[ax] (0,0) -- (0,1.6);
  \node[axl, anchor=west] at (0.06,1.5) {$b_i$};
  \node[axl, anchor=north] at (1.6,-0.02) {frame \#};
  \items{3.2}{1.9}{0/0,0/0,.438/0,.688/0,.25/1,.531/0,0/0,0/0,0/0,.5/1,
                   0/0,.25/0,.344/0,.5/1,.406/0,.312/0,0/0,.25/1,0/0,0/0}
\end{scope}

% ---------------- (3) inference ----------------
% the frozen pruner is a container that holds its output: the frames at their token counts and the text events
\node[model, minimum width=9.4cm, minimum height=1.0cm] (pr) at (7.3,0.7) {};
\node[font=\sffamily\footnotesize, text=modelT, align=center] at (3.4,0.7) {Token\\pruner};
\snowflake{($(pr.north east)+(-0.14,-0.14)$)}{0.075}
\prunedframe{4.95,0.7}{frame1.jpg}{6}
\txtchip{6.15,0.7}
\prunedframe{7.35,0.7}{frame3.jpg}{22}
\prunedframe{8.55,0.7}{frame4.jpg}{12}
\txtchip{9.75,0.7}
\prunedframe{10.95,0.7}{frame5.jpg}{32}
\node[model, minimum width=2.4cm, minimum height=1.0cm] (vlm) at (14.0,0.7) {VLM};
\snowflake{($(vlm.north east)+(-0.14,-0.14)$)}{0.075}
\node[qa] (ans) at (16.4,0.7) {\textbf{A:} Blue};
\draw[arr] (pr.east) -- (vlm.west);
\draw[arr] (vlm.east) -- (ans.west);
\end{tikzpicture}
}
  \caption{\textbf{\name overview.} \ayush{overview description. make the figure ore beautiful.}}
  \label{fig:overview}
\end{figure*}

% ---------------------------------------------------------------------------------------------------------------------
\subsection{Overview and Problem Formulation}
\label{sec:overview}

Frame-level RAG answers a user query $q$ about a video $V$ by passing the top-$K$ most semantically relevant frames to a VLM~\cite{}. Let $f_1, \dots, f_N$ be frames sampled uniformly from $V$, and let $r(f_i\mid q)$ denote the semantic similarity of $f_i$ with respect to $q$ (e.g., CLIP score~\cite{}), normalized across all frames, so that $r(\cdot\mid q)$ approximates the probability that each frame holds the evidence for $q$. Viewed from a mathematical angle, our frame allocation decision is a binary admit or reject, subject to a budget of $K$ frames, so as to maximize the total similarity of the admitted frames. Frame-level RAG is therefore naturally modeled as an integer linear program:
\begin{equation}
  \max_{\mathbf{x} \in \{0,1\}^N} \sum_{i=1}^{N} r(f_i\mid q)\, x_i \qquad \text{s.t.} \quad \sum_{i=1}^{N} x_i \le K.
  \label{eq:topk}
\end{equation}
where $f_i$ is passed to the VLM if $x_i=1$ and dropped if $x_i=0$. We lift this frame-level objective to a multimodal token-level objective, which carefully defines \textit{how} tokens should be scored, and prove that our solution is theoretically more optimal for this objective compared to naive adaptive token allocation baselines (\cref{sec:allocation}). In multi-turn conversations, however, naively re-allocating tokens at every turn is inefficient and breaks prefix caching~\cite{}, a standard feature of serving pipelines. We therefore reformulate our optimization objective to reuse the KV cache across turns. This improves both efficiency and accuracy, and we prove that the cumulative prefill cost stays bounded over long conversations due to our token pruner design (\cref{sec:multiturn}).


% ---------------------------------------------------------------------------------------------------------------------
\subsection{Token Allocation}
\label{sec:allocation}

Moving to token level RAG replaces the binary decision $x_i$ by $b_i$, where $b_i$ represents the number of tokens allocated to frame $f_i$, and the frame budget $K$ by a token budget $B$. Doing so, however, poses two key challenges. First, the objective must state what a frame is worth at every token count. Second, near-duplicate frames that repeat the content of others may be selected, so the relevance of a frame must account for already selected frames (inaccurate? as its account for frames weighted by the tokens allocated to them?). \name addresses the first challenge with a token value function calibrated offline (you can calibrated offline and continously calibrated across examples, then now R can be "updated" which you can define in a simple way), and the second by conditioning relevance on the allocation. We formulate both for frames alone and then extend the formulation to the transcript of the video. (not well written)

\paragraph{Measuring Token Value.}
Intuitively, the token count of a frame acts as its effective resolution. With few tokens, the VLM perceives only the coarse structure of the frame, such as the scene layout and large objects; with more tokens, it also perceives fine detail, such as small objects and on-screen text. How this added detail translates into downstream accuracy, however, is not known in advance, yet it decides how many tokens a frame deserves. We capture it with a \textit{token value function} $v(b)$, the probability (prefer not to say probability as reviewers may get angry) that the VLM reads the evidence of a frame shown with $b$ tokens, and lift~\eqref{eq:topk} to (did not mention how we got v, e.g. we took N video query pairs blah blah)
\begin{equation}
  \max_{\mathbf b}\ \sum_{i}r(f_i\mid q)\,v(b_i)\qquad\text{s.t.}\quad\sum_{i}b_i\le B .
  \label{eq:relevance}
\end{equation}
The objective is thus the probability that the evidence is both in the prompt and readable (dont say probability as its not really literal probability). Maximizing it is the (maybe put "Fig. Y" somewhere thsi is the figure for this offline curve ill add later)
classical problem of optimal search~\cite{koopman1956,stone1975}, with relevance as the prior over the location of the
evidence and tokens as the search effort. Top-$K$ retrieval~\eqref{eq:topk} is the special case $b_i\in\{0,b_{\max}\}$
and $B=Kb_{\max}$, where $b_{\max}$ is the token count of a full frame.

The objective is only as good as its value function. A linear $v$ would reward every token equally and ignore that a
frame saturates, whereas a step at a fixed count would reproduce the all-or-nothing cost of~\eqref{eq:topk}. \name
instead measures $v$ once per VLM, offline. Across a set of calibration videos and queries, it uniformly varies the
number of tokens per frame and fits
\begin{equation} (hmm i change my mind should mention value function FIRST then the objective function)
  v(b)=\frac{b^{\alpha}}{b^{\alpha}+\tau^{\alpha}}
  \label{eq:curve}
\end{equation}
to the probability that the VLM assigns to the correct answer, where $\tau$ is the token count at which a frame reaches
half of its value and $\alpha$ sets how sharply the value rises. (do not say per VLM, just stick to it. We measure V offline, then later sort of and it is updated online and we can sort of define that. Yeah, and kind of fix up the writing, make it write better. No need to say top K retrieval is a special case.

It's kind of obvious. Also never, never say this. Top K retrieval bracket one. That's not how you refer to things.

You say equation. Equation one, equation two. You don't just say brackets one. Hope people know what that means. No one knows what that means.

Maximizing two alone. No, you don't say maximizing equation 2. What is 2? Is it an equation? Is it a function? What is it? Never say stuff like this.

Also, no need to say top K retrieval. No need to mention classical problem of optimal search again stick to the topic measuring token value right we're giving intuition and we're essentially wanting to say okay well that's the intuition uh we need this and we measure this by this offline curve.

Okay, and this is the function. It's updated online. Make and you're putting this into objective function. Objective function looks like this. So there, I explained what v is.

I showed you the equation for v offline as well as online and then I showed you what the objective function is. And done, that's literally it. There's nothing else. Yeah. Very straightforward. Right? No need to mention all this optimal search. What is the relevance to that?

There is no relevance as a prior over the location of no, there is no prior, there is no location. Stop adopting this Bayesian framework. There is nothing Bayesian about. There is no prior.

There is no probabilities here like it's we use probability once like uh try to avoid that term because that's very inaccurate like the reviewer will be angry at you so again so see how I explain very logically follows very logically so that's the measuring token value)

\paragraph{Incorporating Frame Diversity.}
Maximizing~\eqref{eq:relevance} alone, however, spreads the budget too thin. Because the measured curve has $\alpha<1$,
its slope is unbounded at zero tokens, so the optimum gives tokens to every frame with nonzero relevance. Most of these
frames add little, since consecutive frames often show the same content, yet $v$ is measured one frame at a time and
credits every copy in full. Each additional frame is also a potential distractor for the VLM. The root cause is that
$r(f_i\mid q)$ depends on the query alone, whereas a frame tells the VLM little once a similar frame is already shown in
detail. We, therefore, condition the relevance of a frame on the allocation as well:
\begin{equation}
  r(f_i\mid q,\mathbf b)=r(f_i\mid q)-\eta\sum_{j\neq i}s_{ij}\,v(b_j) ,
  \label{eq:cond}
\end{equation}
where $s_{ij}$ measures whether $f_i$ and $f_j$ show the same content and $\eta$ sets how strongly the coverage by
similar frames discounts relevance. Pixel differences are poorly suited for $s_{ij}$, since camera motion changes pixels
without changing content. Temporal distance is also insufficient, since a scene can return much later in the video. We,
therefore, use the cosine $s_{ij}=\max\big(0,\cos(f_i,f_j)\big)$ of the image embeddings that the retriever already
computes. With the conditioned relevance, \name solves
\begin{equation}
  \max_{\mathbf b}\ \sum_{i}r(f_i\mid q,\mathbf b)\,v(b_i)\qquad\text{s.t.}\quad\sum_{i}b_i\le B .
% Ordered pairs (j != i) as in cover_allocate.py, so this eta equals the code's beta.
  \label{eq:p1}
\end{equation}
The objective keeps the form of~\eqref{eq:relevance}, and only the relevance changes. A frame thus loses relevance only
when similar frames are shown with many tokens, so the plan prefers frames with different content. Moreover, a frame
whose conditioned relevance falls to zero receives no tokens at all, so redundancy limits how many frames are shown
without a separate price per frame. A larger $\eta$ yields fewer, more distinct frames, while $\eta=0$
recovers~\eqref{eq:relevance}.

\paragraph{Extending to Transcripts.}
Many questions hinge on what is said rather than shown, so \name also allocates tokens to the transcript of the video.
The extension adds terms of the same form rather than a new mechanism. Let $u_1,\dots,u_M$ be the transcript events,
and index frames and events jointly as items $i\in[N+M]$, with $b_i$ the tokens of item $i$. We abbreviate the query
relevance of item $i$ as $r_i$, which is $r(f_i\mid q)$ for a frame and $\beta\,r(u_m\mid q)$ for an event $u_m$, the
relevance of its text to $q$ scaled by a weight $\beta$. Likewise, $v_i$ is the value function of the type of item $i$:
$v_{\text{frame}}$, given by~\eqref{eq:curve}, or $v_{\text{text}}$, which takes the same form with its own half-value
count and is fitted jointly with $v_{\text{frame}}$, so that both curves are expressed in the same units of answer
probability. To shorten an event, \name uses a prompt compressor that keeps the original words~\cite{llmlingua2}, so
that a shortened event is a subset of the full one rather than a paraphrase. Finally, $s_{ij}=0$ whenever $i$ or $j$ is
an event, so the relevance of an event is not conditioned on the plan. Expanding~\eqref{eq:cond}, the objective
of~\eqref{eq:p1} over all items becomes
\begin{equation}
  J(\mathbf b)=\sum_{i}r_iv_i(b_i)-\eta\sum_{i\neq j}s_{ij}\,v_i(b_i)\,v_j(b_j) .
  \label{eq:joint}
\end{equation}

\paragraph{Iterative Water-Filling.}
The conditioned relevance couples every pair of frames, so maximizing~\eqref{eq:joint} under the budget is a quadratic
knapsack problem with no separable solution. However, with the other items fixed, $J$ is linear in $v_i(b_i)$. Its
coefficient, clipped at zero,
\begin{equation}
  \tilde r_i=\max\Big(0,\ r_i-2\eta\sum_{j\neq i}s_{ij}\,v_j(b_j)\Big) ,
  \label{eq:effective}
\end{equation}
is the residual relevance of item $i$. The factor of two arises because $v_i(b_i)$ enters $J$ twice: it scales the
conditioned relevance of item $i$, and it lowers the conditioned relevance of every item similar to $i$. An item is
therefore charged both for the coverage it receives from similar items and for the coverage it provides to them. The
clipping changes nothing, since an item with a negative coefficient receives no tokens either way. For fixed $\tilde r$,
the problem separates across items, and pricing each token at $\lambda\ge0$ lets every item choose its own token count:
\begin{equation}
  b_i^*(\lambda)=\argmax_{b}\ \tilde r_iv_i(b)-\lambda b ,
  \label{eq:percost}
\end{equation}
where $b$ ranges over the token counts that the pruner supports, including zero. Since the total tokens of the plan
are non-increasing in $\lambda$, a bisection on $\lambda$ finds the price at which the plan meets $B$. It searches
between $\lambda=0$, where every item takes the largest count it can use, and $\max_i\tilde r_i$, where no item is
funded, because $v_i$ never exceeds one and every shown item costs at least one token. We refer to this procedure as
water-filling, with $\lambda$ as the water level.

\begin{proposition}[Everett~\cite{everett1963}]
\label{thm:everett}
For any $\lambda\ge0$, let $B'=\sum_ib_i^*(\lambda)$. Then
\begin{equation}
  \mathbf b^*(\lambda)\in\argmax_{\mathbf b}\Big\{\sum_i\tilde r_iv_i(b_i)\ :\ \sum_ib_i\le B'\Big\} .
\end{equation}
\end{proposition}

\noindent Proposition~\ref{thm:everett} (proof in the supplementary material) yields two implications. First, it
requires neither convexity nor continuity, so each round is exact on the discrete token grid. Second, every funded item
that is not at a bound satisfies $\tilde r_iv_i'(b_i)=\lambda$. All funded items, frames and transcript events alike, thus
earn the same marginal value per token, so tokens flow to whichever item yields more evidence, and an item whose
effective relevance reaches zero receives no tokens.

\name then alternates between the two steps (\cref{alg:iwf}). The first round starts from $\bar v_i=0$, so
that $\tilde r=r$ and the round maximizes~\eqref{eq:joint} with $\eta=0$. Each later round recomputes~\eqref{eq:effective} with
$v_j(b_j)$ replaced by $\bar v_j$, a running average of $v_j(b_j)$ over the previous plans, and water-fills again. The average damps oscillation, since two similar
frames would otherwise trade their tokens back and forth between rounds. \name returns the plan that scores highest
under~\eqref{eq:joint}, and each round costs one water-filling and one product with the similarity matrix. Finally, \name
merges the selected frames and events in time order, each preceded by its timestamp, into one storyboard for the VLM.

\begin{algorithm}[t]
\caption{Iterative Water-Filling Token Allocation}
\label{alg:iwf}
\small
\begin{algorithmic}[1]
\Require relevances $r_i$, value curves $v_i(b)$ for $b\le b_{\max}$ \ayush{too many symbols? fix later}
\Require similarities $s_{ij}$, budget $B$, weight $\eta$, rounds $L$, tolerance $\epsilon$
\Ensure allocation $\mathbf{b}$ with $\sum_i b_i\le B$
\Statex Each update with index $i$ applies to every item $i\in[N+M]$.
\State $\bar v_i\gets 0$ \Comment{coverage of item $i$}
\For{$\ell=1,\dots,L$}
  \State $\tilde r_i\gets\max\!\big(0,\;r_i-2\eta\textstyle\sum_{j\neq i}s_{ij}\bar v_j\big)$ \Comment{residual relevance}
  \State $\lambda_{\mathrm{lo}}\gets 0,\;\lambda_{\mathrm{hi}}\gets\max_i\tilde r_i$ \Comment{price bracket}
  \State $b_i^{(\ell)}\gets 0$ \Comment{feasible at $\lambda_{\mathrm{hi}}$}
  \While{$\lambda_{\mathrm{hi}}-\lambda_{\mathrm{lo}}>\epsilon$} \Comment{bisect on price}
    \State $\lambda\gets(\lambda_{\mathrm{lo}}+\lambda_{\mathrm{hi}})/2$
    \State $b_i\gets\argmax_{b}\;\tilde r_iv_i(b)-\lambda b$ \Comment{best response}
    \If{$\sum_i b_i>B$}
      \State $\lambda_{\mathrm{lo}}\gets\lambda$ \Comment{over budget}
    \Else
      \State $\lambda_{\mathrm{hi}}\gets\lambda$,\; $b_i^{(\ell)}\gets b_i$ \Comment{feasible}
    \EndIf
  \EndWhile
  \State $\bar v_i\gets v_i(b_i^{(\ell)})$ if $\ell=1$, else $\tfrac12\big(\bar v_i+v_i(b_i^{(\ell)})\big)$
\EndFor
\State \Return $\mathbf{b}^{(\ell^\star)}$ with $\ell^\star=\argmax_{\ell}\,J(\mathbf{b}^{(\ell)})$, where
\Statex \hspace{\algorithmicindent}$J(\mathbf{b})=\sum_i r_iv_i(b_i)-\eta\sum_{i\neq j}s_{ij}v_i(b_i)v_j(b_j)$
\end{algorithmic}
\end{algorithm}

\paragraph{Optimality over Prior Allocation Rules.}
Prior allocation rules are feasible points of~\eqref{eq:relevance}: top-$K$ retrieval shows the $K$ most relevant
frames in full, proportional allocation~\cite{luo2025quota} gives frame $i$ the tokens $r_iB$, and rank
tiers~\cite{zhang2025qframe} assign fixed counts by rank. By Proposition~\ref{thm:everett}, on frames the first round of
\cref{alg:iwf} solves~\eqref{eq:relevance} exactly and therefore scores at least as high as each of these rules at
equal spend. Because \cref{alg:iwf} keeps the best plan it visits, its output never scores lower under~\eqref{eq:joint}
than this first round. The gap to proportional allocation, moreover, is structural rather than an artifact of the
token grid:

\begin{proposition}
\label{prop:proportional}
Let $\eta=0$, consider frames only, let the $r_i$ take at least two distinct values, and let every $r_iB$ lie strictly
inside the token range. The proportional plan $b_i=r_iB$ satisfies the optimality condition of~\eqref{eq:relevance}
for all such $r$ if and only if the value curve is logarithmic:
\begin{equation}
  r_i\,v_{\text{frame}}'(r_iB)=\lambda\ \ \forall i \quad\Longleftrightarrow\quad v_{\text{frame}}\ \text{is affine in}\ \log b .
\end{equation}
\end{proposition}

\noindent Proposition~\ref{prop:proportional} (proof in the supplementary material) yields two implications. First,
proportional allocation is optimal only for a value curve that never saturates. Second, since the measured
curve~\eqref{eq:curve} saturates, proportional allocation over-funds the most relevant frames past saturation and
starves the frames behind them, which~\eqref{eq:relevance} funds instead.

% ---------------------------------------------------------------------------------------------------------------------
\subsection{Allocation across Multi-Turn Conversations}
\label{sec:multiturn}

In practice, a user asks several questions $q_1,\dots,q_T$ about one video, and the serving engine keeps the KV cache
of the conversation between turns. Frame-level RAG fits this setting naturally: a later question only appends the
frames it adds, and prefix caching reuses everything before them. Token allocation, however, breaks this reuse, since
solving~\eqref{eq:p1} afresh for each question changes the token counts of frames already in the context, even though
consecutive questions share much of their evidence. An effective multi-turn allocation must therefore (i) reuse the
cached tokens, (ii) adapt relevance to each question, and (iii) bound the prefill of the conversation. Re-allocating
every turn from scratch meets (ii) but neither (i) nor (iii), whereas freezing the first plan meets (i) but not (ii).

\paragraph{Cache-Aware Objective with Attention-Guided Relevance.}
Let $c_i$ be the tokens of frame $i$ already in the cache, with $c_i=0$ before the first turn. Each turn $t$
solves~\eqref{eq:p1} with a turn-specific relevance $r_i^{(t)}$ and prefills only what the cache lacks:
\begin{equation}
  p_i=\max(0,\ b_i-c_i) .
  \label{eq:cost}
\end{equation}
The cache thus changes what a turn costs but not what it chooses. This choice follows KV-cache retrieval and
compression methods, in which the budget bounds the tokens that the model attends to and prefill is a separate
cost~\cite{rekv2025,commitkv2026}. Relevance, however, should change once frames are cached, since the retriever scores
each frame against $q_t$ in isolation, whereas the VLM has already read the cached frames. Before adding new tokens,
\name prefills $q_t$ on top of the cache, which the turn requires anyway, and reads the attention from its tokens to
the cached visual tokens in a few heads, reusing the cached keys as in-context KV-cache retrieval
does~\cite{rekv2025}. Only the attention of the current question is used, since attention under one question misranks
what a later question needs~\cite{cvma2026}. With $a_i$ the share of this attention per cached token of frame $i$,
\name keeps $r_i^{(t)}=r_i$ for frames that are not cached and sets
\begin{equation}
  r_i^{(t)}\propto r_i^{\,1-\gamma}\,a_i^{\,\gamma}\qquad\text{for cached frames } (c_i>0),
  \label{eq:fusion}
\end{equation}
where $\gamma\in[0,1]$ sets the trust in attention, and the proportionality constant keeps the total relevance of the
cached frames equal to their retrieval total.

\begin{proposition}
\label{prop:fusion}
Let $\hat r$ be the retrieval relevance restricted to the cached frames and normalized to sum to one. Among all
distributions $\pi$ on the cached frames,
\begin{equation}
\begin{aligned}
  &\argmin_{\pi}\ (1-\gamma)\,\mathrm{KL}(\pi\,\|\,\hat r)+\gamma\,\mathrm{KL}(\pi\,\|\,a)\\[-2pt]
  &\qquad\propto\ r_i^{\,1-\gamma}a_i^{\,\gamma} .
\end{aligned}
\end{equation}
\end{proposition}

\noindent Proposition~\ref{prop:fusion} (proof in the supplementary material) yields two implications. First,
\eqref{eq:fusion} gives the distribution closest to both signals, from retrieval alone at $\gamma=0$ to
attention alone at $\gamma=1$. Second, because the normalization fixes their total, attention re-ranks the cached frames but
cannot pull budget away from uncached ones, so a question about a new part of the video still sends its tokens there.

\paragraph{Bounded Prefill via Rank-Only Pruning and Append-Only Layout.}
Prefilling only the missing tokens requires two properties. First, the tokens that a frame keeps at a smaller
allocation must be a subset of those it keeps at a larger one. VisionZip~\cite{yang2024visionzip} violates this
property, since it merges the tokens it drops into contextual tokens whose content changes with the allocation. \name
therefore discards the contextual tokens and ranks the dominant tokens of each frame once, independently of the
question, so that a frame at $b$ tokens holds its top $b$ tokens. A turn then prefills exactly the new ranks of each
frame, the $p_i$ tokens of~\eqref{eq:cost}, and updates the cache as
\begin{equation}
  c_i\leftarrow\max(c_i,\ b_i) .
  \label{eq:incr}
\end{equation}

\begin{proposition}
\label{prop:bound}
Let $b_i^{(t)}$ and $p_i^{(t)}$ be the allocation and the prefill of frame $i$ at turn $t$. Over a conversation of $T$
turns,
\begin{equation}
  \sum_{t=1}^{T}p_i^{(t)}=\max_t\,b_i^{(t)}\le b_{\max} .
\end{equation}
\end{proposition}

\noindent Proposition~\ref{prop:bound} (proof in the supplementary material) bounds the prefill of every frame by its
largest allocation, however long the conversation runs, so the whole conversation prefills at most the full tokens of
the frames it ever shows. Re-allocating from scratch, in contrast, costs $\sum_t\sum_ib_i^{(t)}$, which grows with
every turn.

The second property is that the new tokens must not invalidate the cache. Inserting them inside the block of their
frame would invalidate every later token, so \name instead appends them after the conversation in rank order, between
vision start and end tokens. Because the multimodal rotary position embedding of the VLM reserves a position for every
frame, an appended token keeps the temporal and spatial position it has in its frame. Only its place in the causal
order changes.
